% Sharpened theorem-facing refinement of the Section 8--9 ledger.
% Inserted after SECTION9_EXPLICIT_GLOBAL_LEDGER_V3 and before Section 10.

\subsection{Fused endpoint transport and a critical-quarter residual split}
\label{subsec:fused-quarter-refinement-v4}

The preceding ledger is sufficient for the main theorem.  We record here a
stronger form that also shortens the formalization interface.  The first
refinement keeps the endpoint type during the deficit sum and then fuses that
sum with endpoint transport.  The second replaces the auxiliary cubic split
$2^U\le m_0^3$ by the fixed-buffer critical scale $64U2^{U/4}$.  Together they give an
explicit seed exponent of order $\sqrt n(\log n)^{3/2}$.

For all sufficiently large $n$, every $\rho_{ij}$ is smaller than $1/2$.
Put
\begin{equation*}
  B_{ij}:=\frac1{1-\rho_{ij}},
  \qquad
  B^L:=\prod_{i,j}B_{ij}^{\ell_{ij}}.
  \tag{9.43}
\end{equation*}

\begin{lemma}[Type-preserving deficit product]
\label{lem:type-preserving-deficit-product-v4}
For every fixed high-cell support $P$,
\begin{equation*}
  \sum_{h:\,m-h\text{ high}} w(P,m-h)
  \le
  w_{\mathrm{full}}(P)
  \prod_{e\in P}B_{\tau(e)},
  \tag{9.44}
\end{equation*}
where $\tau(e)=(i,j)$ is the endpoint type of $e$.
\end{lemma}

\begin{proof}
For a cell of type $(i,j)$, equations
\eqref{eq:aggregate-deficit-comparison-v3} and
\eqref{eq:three-quarter-charge-v3} give
\[
  n^hR_{m_e,d_e}(h)\le \rho_{ij}^h.
\]
The optional-choice identity may therefore be bounded cell by cell by
\[
  1+\sum_{h\in A_e}\rho_{ij}^h
  \le
  1+\sum_{h\ge1}\rho_{ij}^h
  =\frac1{1-\rho_{ij}}
  =B_{ij}.
\]
Multiplication over the matching support proves the claim.  In contrast with
\eqref{eq:fixed-support-all-deficit-v3}, no factor $\alpha+1$ is introduced.
\end{proof}

Weighted regrouping by the realized endpoint table gives immediately
\begin{equation*}
  \operatorname{BareSkeletonSum}_n
  \le
  \sum_{L\ \mathrm{realized}}W(L)B^L.
  \tag{9.45}
\end{equation*}
Define the fused kernel and its maximal row sum by
\begin{equation*}
  \widetilde Q_{ij}:=B_{ij}Q_{ij},
  \qquad
  \Sigma_n:=\max_{0\le i\le3}\sum_{j=0}^{3}\widetilde Q_{ij}.
  \tag{9.46}
\end{equation*}

\begin{proposition}[Fused deficit--endpoint transport]
\label{prop:fused-deficit-endpoint-transport-v4}
For all sufficiently large $n$,
\begin{equation*}
  \operatorname{BareSkeletonSum}_n
  \le
  \Sigma_n^K\sum_rD(r)
  \le
  \Sigma_n^K\bigl(1+\varepsilon_n^{\mathrm{pd}}\bigr).
  \tag{9.47}
\end{equation*}
\end{proposition}

\begin{proof}
Multiplying \eqref{eq:square-free-endpoint-transport-v3} by $(B^L)^2$ gives
\[
  \bigl(W(L)B^L\bigr)^2
  \le
  \bigl(D(r)A_L\widetilde Q^L\bigr)
  \bigl(D(c)C_L\widetilde Q^L\bigr).
\]
After the arithmetic--geometric mean inequality, sum the first term at fixed
row margin $r$ and discard the column-margin constraint.  The multinomial
theorem gives
\[
  \sum_{L:\,\operatorname{row}(L)=r}A_L\widetilde Q^L
  \le
  \prod_i\left(\sum_j\widetilde Q_{ij}\right)^{r_i}
  \le
  \Sigma_n^{\sum_i r_i}
  \le
  \Sigma_n^K.
\]
The column term is identical because $B$ and $Q$ are symmetric.  Summation over
the two one-sided margins proves the first inequality in (9.47), and
(9.32) proves the second.  Thus the support-cardinality estimate and
the endpoint transport are consumed in one kernel bound.
\end{proof}

The exact phase formula from Section~2 sharpens the deficit scale.  Put
$q:=\log 2$, $L:=\log n$, $\ell:=\log\log n$, and
$C:=1+\log q-q$.  Since
\[
  \alpha=\frac{2(L-\ell+C)}q+b,
  \qquad 0\le b\le1,
\]
and every endpoint minimum satisfies $m_{ij}\ge\alpha-5$, one has
\[
  \left\lfloor\frac{3m_{ij}-1}{4}\right\rfloor
  \ge
  \frac{3\alpha-20}{4}.
\]
Consequently
\[
  q\left\lfloor\frac{3m_{ij}-1}{4}\right\rfloor
  \ge
  \frac32(L-\ell+C)-5q,
\]
and therefore
\begin{equation*}
  2^{\lfloor(3m_{ij}-1)/4\rfloor}
  \ge
  c_0\frac{n^{3/2}}{(\log n)^{3/2}},
  \qquad
  c_0:=\exp\!\left(\frac32C-5q\right)>0.
  \tag{9.48}
\end{equation*}
Since $m_{ij}=O(L)$ uniformly in the phase,
\begin{equation*}
  \rho_{16}
  =O\!\left(\frac{L^{5/2}}{\sqrt n}\right),
  \qquad
  \Sigma_n
  =1+O\!\left(\frac{L^{5/2}}{\sqrt n}\right).
  \tag{9.49}
\end{equation*}
Indeed, $B_{ij}=1+O(\rho_{ij})$, while $Q_{ii}=1$ and the total off-diagonal
mass in each row is $O(\eta_n)=O(L^{3/2}/\sqrt n)$.

Define the sharpened skeleton exponent
\begin{equation*}
  \Gamma_n^{\mathrm{skel},\sharp}
  :=K\log\Sigma_n+
    \log\!\left(1+\varepsilon_n^{\mathrm{pd}}\right).
  \tag{9.50}
\end{equation*}
Since $K=\Theta(n/L)$, equations (9.47)--(9.49) give
\begin{equation*}
  \Gamma_n^{\mathrm{skel},\sharp}
  =O\!\left(\sqrt n\,L^{3/2}\right),
  \qquad
  \operatorname{BareSkeletonSum}_n
  \le e^{\Gamma_n^{\mathrm{skel},\sharp}}.
  \tag{9.51}
\end{equation*}
This improves the earlier coarse deficit contribution
$O(n^{3/4}\log n)$ and removes the artificial factor $\alpha+1$.

We now sharpen the residual split.  Put
\begin{equation*}
  T_U:=64U2^{U/4}.
  \tag{9.52}
\end{equation*}

\begin{proposition}[Critical-quarter residual attachment]
\label{prop:critical-quarter-residual-attachment-v4}
There is an absolute constant $C_\mathrm{q}>0$ such that every feasible
canonical high skeleton satisfies
\begin{equation*}
  \mathcal A(M,j)
  \le
  \begin{cases}
    \exp(C_\mathrm{q}U^2),&m_0\ge T_U,\\[2mm]
    \displaystyle\exp\!\left(32(\log 2)U^2 2^{U/4}\right),&m_0<T_U.
  \end{cases}
  \tag{9.53}
\end{equation*}
The second line includes $m_0=0$.
\end{proposition}

\begin{proof}
Assume first that $m_0\ge T_U$.  The residual degree cap gives
\[
  \theta_{ab}
  =\frac{\mathrm e\,d_ad'_b}{m_0}
  \le \frac{\mathrm e U}{64}2^{-U/4}.
\]
The proof of Lemma~\ref{lem:q-quadratic-v3} now applies with a stronger
endpoint estimate.  For completeness, put
$a_x=g(x)\theta^x/x!$ and $R=\lfloor U/2\rfloor$.  Log-convexity again gives
$\lambda_{ab}\le R(a_3+a_R)$.  The first endpoint satisfies
\[
  \frac{Ra_3}{\theta^2}
  =\frac{2R}{3}\theta
  \le \frac{\mathrm e U^2}{192}2^{-U/4}
  \longrightarrow0.
\]
For the other endpoint, using $g(R)=2^{\binom R2-1}$ for $R\ge3$,
\[
\begin{aligned}
  \log_2\!\left(\frac{Ra_R}{\theta^2}\right)
  \le{}&
  \log_2R+\binom R2-1-\log_2(R!)\\
  &+(R-2)\left(
    \log_2\mathrm e+\log_2U-6-\frac U4
  \right).
\end{aligned}
\]
Here $U\in\{2R,2R+1\}$.  Therefore, for $R\ge3$,
\[
  \binom R2-\frac U4(R-2)\le\frac R2,
  \qquad
  U\le\frac73R.
\]
Moreover, the integral bound for the factorial gives
\[
  \log_2(R!)
  \ge R\log_2R-R\log_2\mathrm e.
\]
Using the elementary estimates
$1<\log_2\mathrm e<3/2$ and
$1<\log_2(7/3)<5/4$, the displayed endpoint logarithm is at most
\[
  7-\frac54R-\log_2R,
\]
which is negative for $R\ge6$.  The finitely many remaining values are
absorbed into one absolute constant.
Thus $q_{ab}\le C_0\theta_{ab}^2$.  The exact identity for
$\sum_{a,b}\theta_{ab}^2$ and the degree cap then give
$\sum_e q_e\le C_1U^2$, so \eqref{eq:q-only-attachment-v3} yields the first
line of (9.53).

If $m_0=0$, the conclusion follows from
\eqref{eq:zero-residual-attachment-v3}.  If $0<m_0<T_U$, the crude estimate
\eqref{eq:complementary-attachment-v3} gives
\[
  \mathcal A(M,j)
  \le2^{Um_0/2}
  \le2^{32U^2 2^{U/4}},
\]
which is the second line.
\end{proof}

Define
\begin{equation*}
  \Gamma_n^{\mathrm{att},\sharp}
  :=
  \max\!\left\{
    C_\mathrm{q}U^2,
    32(\log 2)U^2 2^{U/4}
  \right\}.
  \tag{9.54}
\end{equation*}
The phase relation $2^U=\Theta(n^2/L^2)$ implies
\begin{equation*}
  \Gamma_n^{\mathrm{att},\sharp}
  =O\!\left(\sqrt n\,L^{3/2}\right).
  \tag{9.55}
\end{equation*}

\begin{corollary}[Explicit sharpened seed exponent]
\label{cor:sharpened-seed-exponent-v4}
There is an absolute constant $C_\sharp>0$ such that, uniformly in the phase,
\begin{equation*}
  1
  \le
  \frac{\mathbb E Z^2}{(\mathbb E Z)^2}
  \le
  \exp\!\left(C_\sharp\sqrt n\,(\log n)^{3/2}\right).
  \tag{9.56}
\end{equation*}
Consequently
\begin{equation*}
  \mathbb P(Z>0)
  \ge
  \exp\!\left(-C_\sharp\sqrt n\,(\log n)^{3/2}\right).
  \tag{9.57}
\end{equation*}
\end{corollary}

\begin{proof}
Insert the uniform attachment factor
$e^{\Gamma_n^{\mathrm{att},\sharp}}$ into the exact decomposition
\eqref{eq:exact-attachment-decomposition-v3}, and use (9.51) for the remaining
bare-skeleton sum.  Equations (9.51) and (9.55) give (9.56).  The final claim is
Paley--Zygmund.
\end{proof}

The refinement leaves the main theorem and its coefficient unchanged.  Its
formal advantage is that the high-skeleton calculation is reduced to one
finite kernel row-sum theorem, while the residual calculation is reduced to
the two explicit cases $m_0\ge64U2^{U/4}$ and
$m_0<64U2^{U/4}$.  Neither interface requires
the arbitrary exponent $1/3$ or the factor $\alpha+1$ from the coarser ledger.
